\section{Set Up to Write Together}

The advice of the last section needs a team.
Team members have different comfort with tools.
Some live in the terminal.
Some never used git, and they do not need to.
If we make everyone use one tool, we lose people at both ends.
Thus, we give a ladder, and each author selects a rung:
\begin{enumerate}
\item A reader.
The reader looks at the built PDF and sends comments.
\item A browser writer.
This person writes in a shared Google Doc, with help from claude.ai.
\item An Overleaf writer.
This person edits the LaTeX in the browser and sees the PDF change.
\item A git writer.
This person clones the repository, edits, and pushes.
\item An agent runner.
This person works with an agent, for example Claude Code, on their own
computer.
The agent reads and writes the shared source.
\end{enumerate}
Each rung is a full member of the team.
Only the setup cost is different.

One piece of plumbing connects the ends of the ladder
(Figure~\ref{fig:faces}).
An Overleaf project is also a git remote.
Thus, the same paper has two faces.
Humans see a live editor in the browser.
People and agents see a usual repository.
GitHub holds the source of truth.
A push moves edits to Overleaf in seconds, and a pull collects the
text that the browser writers typed.
Thus, many people and many agents write at the same time, and they
only meet at the merge.

Two habits keep the merges small.
First, each sentence ends its line.
Git then compares texts by the sentence.
Two edits only collide when they touch the same sentence.
Second, the prose lives in many small files, one for each section,
under a thin root of input statements.
Writers then claim a section, not the paper.

Comments follow the same idea.
A comment is a macro in the source.
It shows who spoke and what they said.
One make target lists the open comments.
To resolve a comment, delete it.
We do not use the margin comments of web editors, because those
comments do not go through git.

The last ritual is the schedule.
Before a shared writing session, one owner makes sure that each person
can write to something.
Doc links must arrive.
Project invitations must be accepted.
Push access must be shown with one small commit.
An access problem in a session costs each person an hour.
The same problem, solved one day before, costs one person a minute.

The prompts of the team live in the same repository as the prose.
Thus, each rung runs the same words.
Note: one author and one agent wrote this paper with this method, and
they traded pushes through the bridge above.
Table~\ref{tab:bridge} makes the ladder concrete for this repository.

\begin{table*}[t]
\caption{How to start on this paper, by rung.
The rung 2 writer must get an Overleaf invitation from the first
author.
The rung 3 and 4 writers must get GitHub push access from the same
author.
Each LaTeX rung obeys two habits: start a new line after each full
stop, and keep the prose in many small files under a thin root.}
\label{tab:bridge}
\small
\begin{tabular}{@{}lp{0.42\textwidth}p{0.34\textwidth}@{}}
\toprule
rung & one-time setup & daily cycle\\
\midrule
2: Overleaf browser &
accept the project invitation; open the project; set
{\ttfamily Menu\,$\to$\,Settings\,$\to$\,Main document}
to the paper &
edit the prose in {\ttfamily sec2/}; start a new line after each
full stop; comment with
{\ttfamily \textbackslash note\{you\}\{...\}};
{\ttfamily Cmd-Enter} makes the PDF again\\[0.4em]
3: GitHub &
{\ttfamily git clone github.com/timm/how2rite} &
pull; edit; {\ttfamily make fmt}; commit; push\\[0.4em]
4: the git bridge &
rung 3, plus: in Overleaf, {\ttfamily Account
Settings\,$\to$\,Git Integration\,$\to$\,Generate token} (the
username is {\ttfamily git}, the password is the token); then\newline
{\ttfamily git remote add overleaf}\newline
{\ttfamily ~~https://git.overleaf.com/}\newline
{\ttfamily ~~6a665a138dbd6c9b6ef7d809} &
{\ttfamily git pull overleaf main} (this collects the browser
edits); edit; {\ttfamily make fmt}; build; push to the two remotes:
{\ttfamily git push origin main; git push overleaf main:main}\\
\bottomrule
\end{tabular}
\end{table*}

\begin{pause}{the team's tooling}
Ask: can each team member read and write the files of the paper?
Does each member use a technology where they are productive and
comfortable?
Does the team help everyone, also when system skills are very
different?\\[0.2em] To change the outcome: move people to a different
rung of Table~\ref{tab:bridge} (no rung outranks another), send the
missing invitations or tokens, and name a sync captain who carries
edits between the faces of Figure~\ref{fig:faces}.
\end{pause}
